\chapter[Especificación de Requerimientos de Software]{Especificación de Requerimientos de Software}

\emph{Adaptación basada en IEEE Std 830-1998, Software Requirements
Specifications \citep{ieee_830_1998}.}

\section[Alcances]{Alcances}

OnCuidar es una aplicación móvil para el cuidador principal de un paciente
oncológico pediátrico. Se diferencia de las aplicaciones genéricas de salud en
que su contenido educativo y sus criterios de alerta provienen del equipo de
Enfermería que acompaña el proyecto, en que funciona sin conexión y en que cifra
los datos sensibles en el propio dispositivo.

\subsection*{Funcionalidades dentro del alcance}

\begin{itemize}
    \item Crear una cuenta con los datos del cuidador, del paciente, del centro
    de salud y del contacto de emergencia, e iniciar sesión de forma segura.
    \item Gestionar uno o más pacientes y elegir cuál es el paciente activo.
    \item Registrar signos vitales, síntomas de la escala ESAS y observaciones,
    y clasificar cada registro en los niveles normal, alerta o crítico.
    \item Consultar el historial de registros y exportarlo a PDF y Excel.
    \item Consultar una biblioteca de videos, guías, infografías y documentos,
    disponible sin conexión.
    \item Resolver dudas frecuentes mediante un asistente de orientación y una
    sección de preguntas frecuentes.
    \item Programar recordatorios con notificaciones locales.
\end{itemize}

\subsection*{Funcionalidades fuera del alcance}

\begin{itemize}
    \item Diagnosticar, indicar tratamientos o calcular dosis de medicamentos.
    \item Reemplazar la evaluación del equipo de salud o la atención presencial.
    \item Enviar los registros directamente al hospital o integrarse con su ficha
    clínica electrónica.
    \item Generar respuestas con inteligencia artificial generativa.
    \item Ofrecer una aplicación web o un panel de administración para el equipo
    clínico.
\end{itemize}

\section[Objetivo del Software]{Objetivo del software}

\addtocontents{toc}{\protect\setcounter{tocdepth}{1}}
\subsection{Objetivo general}
\addtocontents{toc}{\protect\setcounter{tocdepth}{2}}

El software gestiona la información del seguimiento domiciliario del paciente
oncológico pediátrico, compuesta por sus signos vitales, síntomas, historial,
recordatorios y contenidos educativos, para apoyar al cuidador en el proceso de
cuidado en el hogar. Con ello logra un registro ordenado y trazable de la
evolución del paciente y una orientación oportuna ante signos de alerta.

\addtocontents{toc}{\protect\setcounter{tocdepth}{1}}
\subsection{Objetivos específicos}
\addtocontents{toc}{\protect\setcounter{tocdepth}{2}}

\begin{enumerate}
    \item Almacenar los antecedentes del cuidador y de sus pacientes, cifrando
    los datos sensibles antes de que salgan del dispositivo.
    \item Registrar la valoración clínica diaria del paciente y clasificar su
    nivel de alerta mediante un motor de reglas clínicas.
    \item Entregar contenido educativo validado y respuestas a dudas frecuentes,
    disponibles incluso sin conexión.
    \item Generar informes del historial en PDF y Excel para presentarlos en los
    controles clínicos.
    \item Recordar al cuidador las tareas de cuidado mediante notificaciones
    locales.
\end{enumerate}

\section[Descripción Global del Producto]{Descripción global del producto}

\subsection[Interfaz de Usuario]{Interfaz de usuario}

La interfaz está pensada para personas sin formación técnica ni clínica que
pueden usarla en momentos de estrés. Por eso emplea textos breves, botones
amplios y un lenguaje cotidiano. Sigue las pautas de Material~3 con una paleta
cálida basada en el dorado del lazo del cáncer infantil y la tipografía Nunito.

Todas las pantallas comparten una cabecera con degradado dorado y una barra de
navegación inferior con cinco accesos: Inicio, Chat, Registro, Aprende y
Perfil. El nivel de alerta se comunica siempre con el mismo código de colores,
verde para normal, amarillo para alerta y rojo para crítico, acompañado de un
ícono y un texto para no depender solo del color.

\subsection[Interfaz de Hardware]{Interfaz de hardware}

La aplicación se ejecuta en teléfonos Android e iOS compatibles con Flutter y no
requiere sensores ni periféricos especiales. Utiliza el almacenamiento interno
para la persistencia local de datos, la caché del contenido educativo y los
archivos exportados. Solicita permisos de notificaciones, de alarmas exactas y
de ejecución al reiniciar el dispositivo, para que los recordatorios se
disparen a la hora programada.

\subsection[Interfaz de Software]{Interfaz software}

La Tabla~\ref{tab:interfaz-software} detalla los productos de software con los
que se relaciona el sistema.

\begingroup
\small
\tablazebra
\begin{xltabular}{\textwidth}{P{1.25}P{0.85}P{0.80}P{1.10}}
    \caption{Productos de software con los que se relaciona el sistema.}
    \label{tab:interfaz-software}\\
    \toprule
    \filaencabezado \encab{Nombre} & \encab{Abreviación} & \encab{Versión} & \encab{Fuente} \\
    \midrule
    \endfirsthead
    \rowcolor{white}\multicolumn{4}{@{}l}{\footnotesize\itshape Tabla \thetable{} (continuación)} \\
    \toprule
    \filaencabezado \encab{Nombre} & \encab{Abreviación} & \encab{Versión} & \encab{Fuente} \\
    \midrule
    \endhead
    \midrule
    \rowcolor{white}\multicolumn{4}{r@{}}{\footnotesize\itshape Continúa en la página siguiente} \\
    \endfoot
    \bottomrule
    \endlastfoot
    Flutter SDK (Dart) & Flutter & Dart 3.12 & flutter.dev \\
    http & --- & 1.6.0 & pub.dev \\
    flutter\_riverpod & Riverpod & 3.4.2 & pub.dev \\
    go\_router & --- & 17.5.0 & pub.dev \\
    cryptography & --- & 2.9.0 & pub.dev \\
    flutter\_secure\_storage & --- & 9.2.4 & pub.dev \\
    flutter\_local\_notifications & --- & 22.3.1 & pub.dev \\
    flutter\_cache\_manager & --- & 3.4.4 & pub.dev \\
    connectivity\_plus & --- & 7.3.1 & pub.dev \\
    pdf & --- & 3.12.0 & pub.dev \\
    excel & --- & 4.0.6 & pub.dev \\
    video\_player & --- & 2.14.0 & pub.dev \\
    url\_launcher & --- & 6.3.2 & pub.dev \\
    share\_plus & --- & 12.0.2 & pub.dev \\
    Node.js & Node & 22 & nodejs.org \\
    Express & --- & 5.3.0 & npmjs.com \\
    PostgreSQL & Postgres & 15 & postgresql.org \\
    pg (cliente de PostgreSQL) & --- & 8.23.1 & npmjs.com \\
    bcryptjs & --- & 2.4.3 & npmjs.com \\
    jsonwebtoken & --- & 9.0.3 & npmjs.com \\
    nodemailer & --- & 6.10.1 & npmjs.com \\
    helmet & --- & 8.3.0 & npmjs.com \\
    cors & --- & 2.8.6 & npmjs.com \\
    Jest & --- & 29.7.0 & npmjs.com \\
    supertest & --- & 7.3.1 & npmjs.com \\
    Docker Compose y Podman & --- & --- & docker.com, podman.io \\ % REVISAR: versión instalada en el servidor
\end{xltabular}
\endgroup

\addtocontents{toc}{\protect\setcounter{tocdepth}{1}}
\subsection{Interfaces de comunicación}
\addtocontents{toc}{\protect\setcounter{tocdepth}{2}}

La aplicación se comunica con la API REST propia mediante HTTP, con mensajes en
formato JSON y autenticación por token JWT en cada solicitud. En el entorno del
servidor de la universidad se prevé servir la API sobre HTTPS. El contenido
educativo (videos, infografías y guías) lo sirve el mismo backend desde la ruta
/archivos y la aplicación lo descarga para su uso sin conexión. El
servidor envía los correos de verificación y recuperación mediante SMTP. La
aplicación abre llamadas telefónicas con el esquema tel: y direcciones
en Google Maps mediante enlaces web. No se definen protocolos propios.

\section[Requerimientos Específicos]{Requerimientos específicos}

\subsection[Requerimientos Funcionales del Sistema]{Requerimientos funcionales del sistema}

Los requerimientos funcionales se derivan de las historias de usuario del
backlog y se agrupan por módulo. La columna HU indica la historia de usuario que
da origen a cada requerimiento, lo que permite seguir su trazabilidad hasta las
pruebas.

\begingroup
\small
\tablazebra
\begin{xltabular}{\textwidth}{P{0.42}P{0.95}P{2.33}P{0.55}P{0.75}}
    \caption{Requerimientos funcionales del sistema.}
    \label{tab:requerimientos-funcionales}\\
    \toprule
    \filaencabezado \encab{ID} & \encab{Nombre} & \encab{Descripción} & \encab{HU} & \encab{Prioridad} \\
    \midrule
    \endfirsthead
    \rowcolor{white}\multicolumn{5}{@{}l}{\footnotesize\itshape Tabla \thetable{} (continuación)} \\
    \toprule
    \filaencabezado \encab{ID} & \encab{Nombre} & \encab{Descripción} & \encab{HU} & \encab{Prioridad} \\
    \midrule
    \endhead
    \midrule
    \rowcolor{white}\multicolumn{5}{r@{}}{\footnotesize\itshape Continúa en la página siguiente} \\
    \endfoot
    \bottomrule
    \endlastfoot
    \filagrupo{5}{Acceso y seguridad}
    RF-01 & Iniciar la aplicación &
    Mostrar una pantalla de carga y dirigir al panel principal si existe una
    sesión activa, o a la bienvenida si no la hay. & HU-01 & Crítica \\
    RF-02 & Crear cuenta &
    Registrar la cuenta con correo y contraseña junto con los datos del
    cuidador, del paciente, del centro de salud y del contacto de emergencia,
    validando los campos obligatorios y sus formatos. & HU-02 & Crítica \\
    RF-03 & Iniciar y cerrar sesión &
    Autenticar al cuidador, informar con claridad las credenciales inválidas,
    mantener la sesión abierta entre usos y cerrarla limpiando los datos
    locales. & HU-03 & Alta \\
    RF-04 & Recuperar acceso &
    Enviar al correo principal o al de respaldo un enlace de restablecimiento de
    contraseña de un solo uso, con una hora de vigencia. & HU-03 & Alta \\
    RF-05 & Cifrar datos sensibles &
    Cifrar en el dispositivo los datos personales y clínicos antes de enviarlos al
    servidor, obtener la clave de cada usuario desde el servidor que la custodia y
    bloquear el acceso a la clave al cerrar sesión. & HU-04 & Crítica \\
    RF-29 & Verificar correos por enlace &
    Enviar un enlace de verificación de un solo uso, con una hora de vigencia, al
    correo principal al registrarse o cambiarlo y al correo de respaldo; mostrar
    si el correo está verificado o pendiente, sin bloquear el uso de la
    aplicación, y permitir reenviar el enlace una vez por minuto como máximo. & HU-05 & Media \\
    \filagrupo{5}{Perfil y pacientes}
    RF-06 & Gestionar perfil del cuidador &
    Mostrar y editar los datos del cuidador, exigir la contraseña actual para
    cambiar un correo e indicar cuándo el nuevo correo está pendiente de
    verificación. & HU-05 & Alta \\
    RF-07 & Gestionar pacientes &
    Crear, editar, archivar y restaurar los pacientes del cuidador, con
    confirmación antes de archivar. & HU-06 & Crítica \\
    RF-08 & Seleccionar paciente activo &
    Elegir el paciente activo desde una lista, mostrar su nombre en el panel
    principal y conservar la selección al navegar o cerrar la aplicación. & HU-07 & Crítica \\
    RF-09 & Contactar al centro de salud &
    Iniciar una llamada al centro de salud del paciente y abrir su dirección en
    la aplicación de mapas. & HU-06 & Media \\
    RF-10 & Separar datos por paciente &
    Asociar registros y recordatorios al paciente activo
    y actualizar todas las vistas al cambiar de paciente. & HU-24 & Alta \\
    \filagrupo{5}{Registro clínico e historial}
    RF-11 & Registrar valoración clínica &
    Registrar hasta cuatro signos vitales opcionales (temperatura, frecuencia
    cardíaca, saturación de oxígeno y frecuencia respiratoria), la intensidad de
    los síntomas ESAS de 0 a 10 (sin síntoma, leve, moderado, severo e insoportable) y
    observaciones, rechazando valores fuera de
    rango físico. & HU-09 & Crítica \\
    RF-12 & Controlar tope diario &
    Distinguir registros programados y extra, y bloquear los programados cuando
    el paciente alcanza su tope diario, fijado en tres por defecto. & HU-09 & Crítica \\
    RF-13 & Clasificar nivel de alerta &
    Clasificar cada registro como normal, alerta o crítico según umbrales de
    signos vitales e intensidad de síntomas, y mostrar el motivo junto con la
    recomendación correspondiente. & HU-09 & Crítica \\
    RF-14 & Consultar historial &
    Listar en orden cronológico los registros del paciente activo, filtrarlos
    por fecha o nivel de alerta y mostrar el detalle de cada uno. & HU-10 & Alta \\
    RF-15 & Editar y eliminar registros &
    Modificar o eliminar los registros del día, con confirmación previa a la
    eliminación. & HU-11 & Media \\
    RF-16 & Exportar historial a PDF &
    Generar un informe PDF con los datos del paciente y sus registros en un rango
    de fechas, y permitir guardarlo o compartirlo. & HU-12 & Alta \\
    RF-17 & Exportar historial a Excel &
    Generar una planilla Excel con fecha, signos, síntomas, nivel y
    observaciones de cada registro en un rango de fechas, y permitir guardarla o
    compartirla. & HU-13 & Media \\
    RF-18 & Mostrar panel principal &
    Presentar el paciente activo, la fecha del último registro, un resumen de
    signos vitales, los registros del día, el nivel de alerta vigente y accesos
    directos a los módulos. & HU-08 & Alta \\
    \filagrupo{5}{Orientación}
    RF-19 & Consultar asistente &
    Responder por chat la consulta escrita del cuidador con la pregunta frecuente
    que mejor coincide por palabras clave, o con una respuesta de respaldo si no
    hay coincidencia. & HU-14 & Media \\
    RF-20 & Gestionar conversaciones &
    Guardar las conversaciones con un título automático, listarlas por última
    actividad y permitir continuarlas, renombrarlas y eliminarlas. & HU-15 & Baja \\
    RF-21 & Consultar preguntas frecuentes &
    Mostrar las preguntas frecuentes por categoría, con búsqueda por palabra
    clave que ignora tildes. & HU-16 & Baja \\
    \filagrupo{5}{Recordatorios}
    RF-22 & Gestionar recordatorios &
    Crear, editar, activar o desactivar y eliminar recordatorios de medicamento, medición,
    cita u otro tipo, únicos o recurrentes. & HU-17 & Alta \\
    RF-23 & Notificar recordatorios &
    Emitir una notificación local a la hora de cada recordatorio activo,
    cancelarlas al cerrar sesión y reprogramarlas al volver a iniciarla. & HU-18 & Media \\
    \filagrupo{5}{Educación}
    RF-24 & Consultar biblioteca &
    Mostrar el catálogo de videos, guías, infografías y documentos, con filtro
    por categoría, búsqueda por título, favoritos y estado de descarga. & HU-19 & Media \\
    RF-25 & Descargar contenido &
    Descargar el contenido educativo en segundo plano tras el primer ingreso,
    reintentar ante fallas de red y mostrar el progreso. & HU-20 & Alta \\
    RF-26 & Reproducir videos &
    Reproducir los videos educativos en horizontal y a pantalla inmersiva, con
    controles para reproducir, pausar, avanzar y retroceder 10 segundos, una barra
    de progreso y la reanudación desde donde quedó. & HU-21 & Media \\
    RF-27 & Ver guías e infografías &
    Mostrar las guías como recomendaciones en secciones con título y las
    infografías ampliadas sobre la pantalla, con zoom de hasta cinco veces. & HU-22 & Media \\
    \filagrupo{5}{Funcionamiento sin conexión}
    RF-28 & Operar sin conexión &
    Crear y consultar registros sin conexión, guardar las escrituras en una cola
    persistente que se reintenta al recuperar la red y avisar al cuidador cuando
    no hay conexión. & HU-27 & Alta \\
\end{xltabular}
\endgroup

\subsection[Interfaces Externas de Entrada]{Interfaces externas de entrada}

La Tabla~\ref{tab:entradas} agrupa los datos que ingresan al sistema.

\begingroup
\small
\tablazebra
\begin{xltabular}{\textwidth}{P{0.45}P{0.75}P{1.80}}
    \caption{Interfaces externas de entrada.}
    \label{tab:entradas}\\
    \toprule
    \filaencabezado \encab{Identificador} & \encab{Nombre del ítem} & \encab{Detalle de datos contenidos en el ítem} \\
    \midrule
    \endfirsthead
    \rowcolor{white}\multicolumn{3}{@{}l}{\footnotesize\itshape Tabla \thetable{} (continuación)} \\
    \toprule
    \filaencabezado \encab{Identificador} & \encab{Nombre del ítem} & \encab{Detalle de datos contenidos en el ítem} \\
    \midrule
    \endhead
    \midrule
    \rowcolor{white}\multicolumn{3}{r@{}}{\footnotesize\itshape Continúa en la página siguiente} \\
    \endfoot
    \bottomrule
    \endlastfoot
    DE\_01 & Datos de la cuenta & Correo electrónico, correo de respaldo,
    contraseña y confirmación de contraseña. \\
    DE\_02 & Datos del cuidador & Nombre completo, relación con el paciente,
    teléfono y dirección. \\
    DE\_03 & Datos del paciente & Nombre, RUT, edad, diagnóstico, fase de
    tratamiento y tope diario de registros. \\
    DE\_04 & Centro de salud y emergencia & Nombre, dirección y teléfono del
    centro de salud, nombre y teléfono del contacto de emergencia. \\
    DE\_05 & Credenciales de acceso & Correo electrónico y contraseña. \\
    DE\_06 & Valoración clínica & Tipo de registro, temperatura, frecuencia
    cardíaca, saturación de oxígeno, frecuencia respiratoria, síntomas con su
    intensidad y observaciones. \\
    DE\_07 & Consulta al asistente & Texto libre de la duda del cuidador. \\
    DE\_08 & Criterios de búsqueda & Palabras clave, categoría y rango de
    fechas. \\
    DE\_09 & Recordatorio & Título, descripción, tipo, fecha, hora, días de
    repetición y recurrencia. \\
\end{xltabular}
\endgroup

\subsection[Interfaces Externas de Salida]{Interfaces externas de salida}

La Tabla~\ref{tab:salidas} detalla las salidas del sistema y su medio.

\begingroup
\small
\tablazebra
\begin{xltabular}{\textwidth}{P{0.55}P{0.90}P{1.95}P{0.60}}
    \caption{Interfaces externas de salida.}
    \label{tab:salidas}\\
    \toprule
    \filaencabezado \encab{Identificador} & \encab{Nombre del ítem} & \encab{Detalle de datos contenidos en el ítem} & \encab{Medio} \\
    \midrule
    \endfirsthead
    \rowcolor{white}\multicolumn{4}{@{}l}{\footnotesize\itshape Tabla \thetable{} (continuación)} \\
    \toprule
    \filaencabezado \encab{Identificador} & \encab{Nombre del ítem} & \encab{Detalle de datos contenidos en el ítem} & \encab{Medio} \\
    \midrule
    \endhead
    \midrule
    \rowcolor{white}\multicolumn{4}{r@{}}{\footnotesize\itshape Continúa en la página siguiente} \\
    \endfoot
    \bottomrule
    \endlastfoot
    IS\_01 & Panel principal & Paciente activo, último registro, signos
    vitales recientes, registros del día y nivel de alerta vigente. & Pantalla \\
    IS\_02 & Resultado de la valoración & Nivel de alerta, motivos detectados y
    recomendación. & Pantalla \\
    IS\_03 & Historial & Registros ordenados por fecha con sus signos, síntomas,
    observaciones y nivel de alerta. & Pantalla \\
    IS\_04 & Informe de seguimiento & Datos del paciente, del cuidador, del
    centro de salud y del contacto de emergencia, resumen y detalle de
    registros. & Archivo PDF \\
    IS\_05 & Planilla de registros & Fecha, signos vitales, síntomas, nivel de
    alerta y observaciones de cada registro. & Archivo XLSX \\
    IS\_06 & Respuesta del asistente & Respuesta a la consulta y acceso a las
    preguntas frecuentes. & Pantalla \\
    IS\_07 & Recordatorio & Título y descripción del recordatorio a la hora
    programada. & Notificación \\
    IS\_08 & Aviso de conexión & Estado sin conexión de la aplicación. & Pantalla \\
\end{xltabular}
\endgroup

\subsection[Atributos del Producto]{Atributos del producto}

Los atributos de calidad se definen con base en el modelo ISO/IEC~9126 y se
detallan en la Tabla~\ref{tab:atributos}.

\begingroup
\small
\tablazebra
\begin{xltabular}{\textwidth}{P{0.40}P{0.80}P{1.40}P{1.40}}
    \caption{Atributos del producto.}
    \label{tab:atributos}\\
    \toprule
    \filaencabezado \encab{ID} & \encab{Atributo} & \encab{Descripción} & \encab{Criterio de aceptación} \\
    \midrule
    \endfirsthead
    \rowcolor{white}\multicolumn{4}{@{}l}{\footnotesize\itshape Tabla \thetable{} (continuación)} \\
    \toprule
    \filaencabezado \encab{ID} & \encab{Atributo} & \encab{Descripción} & \encab{Criterio de aceptación} \\
    \midrule
    \endhead
    \midrule
    \rowcolor{white}\multicolumn{4}{r@{}}{\footnotesize\itshape Continúa en la página siguiente} \\
    \endfoot
    \bottomrule
    \endlastfoot
    RNF-01 & Usabilidad y operabilidad &
    La interfaz usa textos breves, lenguaje cotidiano y botones amplios, y cada
    mensaje de error indica qué ocurrió y cómo corregirlo. &
    El 100~\% de los mensajes de validación señala el campo afectado y la acción
    esperada. \\
    RNF-02 & Usabilidad y navegación &
    Los módulos principales se alcanzan desde una barra inferior permanente y
    desde los accesos del panel principal. &
    Cada módulo principal es accesible desde el panel o la barra inferior sin
    pasar por pantallas intermedias. \\
    RNF-03 & Funcionalidad y seguridad &
    El acceso requiere correo y contraseña, los datos sensibles se cifran con
    AES-256-GCM en el dispositivo y cada cuidador accede solo a sus propios
    pacientes. &
    El servidor rechaza con error 400 la escritura de campos clínicos en claro y
    filtra todas las consultas por el usuario autenticado. \\
    RNF-04 & Fiabilidad y disponibilidad &
    La aplicación permite consultar el historial, el contenido descargado y las
    preguntas frecuentes sin conexión. &
    Sin red, las funciones de consulta siguen operativas y los registros nuevos
    se guardan en una cola local que se envía al volver la conexión, sin duplicar
    datos porque las escrituras son idempotentes. \\
    RNF-05 & Eficiencia y tiempo de respuesta &
    Las pantallas principales cargan con rapidez para no entorpecer el uso en
    situaciones de estrés. &
    El panel, el registro y el historial cargan en menos de 3~segundos con una
    conexión móvil estándar. \\
    RNF-06 & Mantenibilidad &
    El contenido educativo y las preguntas frecuentes se mantienen separados de
    la lógica de la aplicación. &
    El contenido se actualiza sin modificar las pantallas ni la lógica de la
    aplicación. \\
    RNF-07 & Portabilidad &
    La aplicación se ejecuta en Android e iOS desde una misma base de código. &
    La aplicación compila y funciona en los dispositivos de prueba de ambas
    plataformas. \\
    RNF-08 & Pertinencia clínica &
    La aplicación no diagnostica, no indica tratamientos ni calcula dosis. &
    Toda orientación se limita a contenido validado y, ante riesgo, recomienda
    contactar al equipo de salud. \\
\end{xltabular}
\endgroup
